\documentclass[12pt,a4paper]{article}
\usepackage[utf8]{vietnam}
\usepackage{amsmath,amssymb}
\usepackage[top=1.8cm,bottom=2cm,left=2cm,right=2cm]{geometry}
\usepackage{fancyhdr}

\pagestyle{fancy}
\fancyhf{}
\rhead{\textit{Trang \thepage}}
\lhead{\textbf{KỲ THI THỬ TỐT NGHIỆP THPT 2027 — MÔN TOÁN}}

\begin{document}

\begin{center}
\textbf{\large BỘ GIÁO DỤC VÀ ĐÀO TẠO} \hfill \textbf{\large ĐỀ KHẢO SÁT CHẤT LƯỢNG NĂM 2027} \\
\textbf{\large ĐỀ THI CHÍNH THỨC} \hfill \textbf{Môn: TOÁN (Chương trình GDPT 2018)} \\
\textit{(Đề thi có 22 câu hỏi)} \hfill \textit{Thời gian làm bài: 90 phút (không kể phát đề)}
\end{center}
\hrule
\vspace{0.4cm}

\textbf{\large PHẦN I. Thí sinh trả lời từ câu 1 đến câu 12. Mỗi câu hỏi thí sinh chỉ chọn một phương án.}

\vspace{0.2cm}
\textbf{Câu 1.} Cho hàm số $y = f(x)$ xác định và liên tục trên $\mathbb{R}$, có bảng xét dấu đạo hàm $f'(x)$ đổi dấu từ dương sang âm tại $x = -1$ và từ âm sang dương tại $x = 2$. Hàm số đã cho đồng biến trên khoảng nào dưới đây? \\
A. $(-1; 2)$ \hfill \textbf{B.} $(2; +\infty)$ \hfill C. $(-\infty; 2)$ \hfill D. $(-1; +\infty)$

\vspace{0.2cm}
\textbf{Câu 2.} Tập xác định của hàm số $y = \log_3(2x - 4)$ là \\
A. $[2; +\infty)$ \hfill B. $(-\infty; 2)$ \hfill \textbf{C.} $(2; +\infty)$ \hfill D. $\mathbb{R} \setminus \{2\}$

\vspace{0.2cm}
\textbf{Câu 3.} Cho $\int_1^3 f(x)\,\mathrm{d}x = 4$ và $\int_1^3 g(x)\,\mathrm{d}x = -1$. Giá trị của $\int_1^3 [2f(x) + 3g(x)]\,\mathrm{d}x$ bằng \\
\textbf{A.} $5$ \hfill B. $11$ \hfill C. $7$ \hfill D. $1$

\vspace{0.2cm}
\textbf{Câu 4.} Trong không gian $Oxyz$, cho mặt cầu $(S): (x - 1)^2 + (y + 2)^2 + (z - 3)^2 = 25$. Tọa độ tâm $I$ và bán kính $R$ là \\
\textbf{A.} $I(1; -2; 3), R = 5$ \hfill B. $I(-1; 2; -3), R = 5$ \hfill C. $I(1; -2; 3), R = 25$ \hfill D. $I(-1; 2; -3), R = 25$

\vspace{0.2cm}
\textbf{Câu 5.} Cho cấp số nhân $(u_n)$ với $u_1 = 3$ và công bội $q = 2$. Giá trị của số hạng $u_4$ bằng \\
\textbf{A.} $24$ \hfill B. $48$ \hfill C. $18$ \hfill D. $12$

\vspace{0.2cm}
\textbf{Câu 6.} Đường tiệm cận ngang của đồ thị hàm số $y = \dfrac{3x - 1}{x + 2}$ có phương trình là \\
\textbf{A.} $y = 3$ \hfill B. $x = -2$ \hfill C. $y = -\dfrac{1}{2}$ \hfill D. $x = 3$

\vspace{0.2cm}
\textbf{Câu 7.} Trong không gian $Oxyz$, cho hai điểm $A(1; 2; -1)$ và $B(3; 0; 1)$. Tọa độ trung điểm $M$ của đoạn thẳng $AB$ là \\
\textbf{A.} $(2; 1; 0)$ \hfill B. $(4; 2; 0)$ \hfill C. $(1; -1; 1)$ \hfill D. $(2; -1; 1)$

\vspace{0.2cm}
\textbf{Câu 8.} Nghiệm của phương trình $2^{x + 1} = 16$ là \\
\textbf{A.} $x = 3$ \hfill B. $x = 4$ \hfill C. $x = 2$ \hfill D. $x = 5$

\vspace{0.2cm}
\textbf{Câu 9.} Thể tích của khối chóp có diện tích đáy $B = 9\text{ cm}^2$ và chiều cao $h = 4\text{ cm}$ là \\
\textbf{A.} $12\text{ cm}^3$ \hfill B. $36\text{ cm}^3$ \hfill C. $18\text{ cm}^3$ \hfill D. $24\text{ cm}^3$

\vspace{0.2cm}
\textbf{Câu 10.} Cho khối lăng trụ đứng $ABC.A'B'C'$ có đáy $ABC$ vuông tại $A$, $AB = 3, AC = 4$ và $AA' = 5$. Thể tích khối lăng trụ bằng \\
\textbf{A.} $30$ \hfill B. $10$ \hfill C. $60$ \hfill D. $20$

\vspace{0.2cm}
\textbf{Câu 11.} Trong không gian $Oxyz$, vectơ nào dưới đây là một vectơ pháp tuyến của $(P): 2x - 3y + z - 7 = 0$? \\
\textbf{A.} $\vec{n} = (2; -3; 1)$ \hfill B. $\vec{n} = (2; 3; 1)$ \hfill C. $\vec{n} = (-2; -3; 1)$ \hfill D. $\vec{n} = (2; -3; -7)$

\vspace{0.2cm}
\textbf{Câu 12.} Đạo hàm của hàm số $y = \mathrm{e}^{2x - 1}$ là \\
\textbf{A.} $y' = 2\mathrm{e}^{2x - 1}$ \hfill B. $y' = \mathrm{e}^{2x - 1}$ \hfill C. $y' = \dfrac{1}{2}\mathrm{e}^{2x - 1}$ \hfill D. $y' = (2x - 1)\mathrm{e}^{2x - 2}$

\vspace{0.4cm}
\textbf{\large PHẦN II. Thí sinh trả lời từ câu 1 đến câu 4. Trong mỗi ý a), b), c), d) ở mỗi câu, thí sinh chọn đúng hoặc sai.}

\vspace{0.2cm}
\textbf{Câu 1.} Cho hàm số $f(x) = x^3 - 3x^2 + 2$ có đồ thị $(C)$. \\
a) Hàm số đồng biến trên các khoảng $(-\infty; 0)$ và $(2; +\infty)$. \\
b) Điểm cực đại của đồ thị hàm số $(C)$ là $A(0; 2)$. \\
c) Giá trị nhỏ nhất của hàm số trên đoạn $[0; 3]$ bằng $-2$. \\
d) Tiếp tuyến của $(C)$ tại điểm có hoành độ $x = 1$ song song với $y = -3x + 5$.

\vspace{0.2cm}
\textbf{Câu 2.} Trong không gian $Oxyz$, cho $A(1; 0; 2)$, $B(2; -1; 3)$, $C(0; 2; 1)$ và $(P): x + y + z - 6 = 0$. \\
a) Vectơ $\vec{AB} = (1; -1; 1)$. \\
b) Điểm $M$ thuộc đoạn thẳng $AB$ sao cho $AM = 2MB$ có tọa độ $\left(\dfrac{5}{3}; -\dfrac{2}{3}; \dfrac{8}{3}\right)$. \\
c) Mặt phẳng đi qua $A$ và song song với $(P)$ có phương trình $x + y + z - 3 = 0$. \\
d) Tam giác $ABC$ là tam giác đều.

\vspace{0.2cm}
\textbf{Câu 3.} Một hộp chứa 6 quả cầu đỏ và 4 quả cầu xanh. Lấy ngẫu nhiên đồng thời 3 quả cầu từ hộp. \\
a) Số phần tử của không gian mẫu là $n(\Omega) = 120$. \\
b) Xác suất để lấy được 3 quả cầu cùng màu đỏ là $\dfrac{1}{6}$. \\
c) Xác suất để lấy được ít nhất một quả cầu màu xanh là $\dfrac{5}{6}$. \\
d) Xác suất để lấy được cả quả cầu đỏ lẫn quả cầu xanh bằng $\dfrac{1}{3}$.

\vspace{0.2cm}
\textbf{Câu 4.} Cho hình chóp $S.ABCD$ có đáy là hình vuông cạnh $a$, $SA \perp (ABCD)$ và $SA = a\sqrt{2}$. \\
a) $BC \perp (SAB)$. \\
b) Góc giữa đường thẳng $SC$ và $(ABCD)$ bằng $45^\circ$. \\
c) Khoảng cách từ điểm $A$ đến mặt phẳng $(SBC)$ bằng $\dfrac{a\sqrt{6}}{3}$. \\
d) Thể tích của khối chóp $S.ABCD$ bằng $\dfrac{a^3\sqrt{2}}{6}$.

\vspace{0.4cm}
\textbf{\large PHẦN III. Thí sinh trả lời từ câu 1 đến câu 6. Điền kết quả vào ô số.}

\vspace{0.2cm}
\textbf{Câu 1.} Một vật chuyển động theo quy luật $s(t) = -t^3 + 9t^2 + t$ ($s$ tính bằng mét, $t$ tính bằng giây, $0 \le t \le 5$). Vận tốc tức thời lớn nhất của vật trong khoảng thời gian đó đạt được tại thời điểm $t$ bằng bao nhiêu giây?

\vspace{0.2cm}
\textbf{Câu 2.} Một doanh nghiệp sản xuất một loại sản phẩm với hàm chi phí $C(x) = x^2 + 40x + 1000$ (nghìn đồng) và doanh thu $R(x) = 160x$ (nghìn đồng), trong đó $x > 0$. Lợi nhuận lớn nhất mà doanh nghiệp có thể đạt được là bao nhiêu nghìn đồng?

\vspace{0.2cm}
\textbf{Câu 3.} Một người muốn làm một bể cá bằng kính không có nắp dạng hình hộp chữ nhật có đáy là hình vuông, thể tích $0{,}5\text{ m}^3$. Hỏi diện tích kính tối thiểu cần dùng là bao nhiêu mét vuông?

\vspace{0.2cm}
\textbf{Câu 4.} Cho hàm số $f(x)$ liên tục trên $\mathbb{R}$ và thỏa mãn $\int_0^3 f(x)\,\mathrm{d}x = 6$. Tính giá trị của $I = \int_0^1 x \cdot f(3x^2)\,\mathrm{d}x$.

\vspace{0.2cm}
\textbf{Câu 5.} Trong không gian $Oxyz$, cho mặt phẳng $(P): 2x - y + 2z - 9 = 0$ và điểm $A(1; 3; -2)$. Tính khoảng cách từ điểm $A$ đến mặt phẳng $(P)$. (Làm tròn đến 2 chữ số thập phân).

\vspace{0.2cm}
\textbf{Câu 6.} Biết phương trình $\log_2^2 x - 3\log_2 x + 2 = 0$ có hai nghiệm $x_1, x_2$. Tính tích $x_1 \cdot x_2$.

\begin{center}
\textbf{--- HẾT ---}
\end{center}

\end{document}
